\begin{abstract}
Difference-in-Differences (DID) confounds treatment effects with differential untreated outcome growth when parallel trends fails. This thesis evaluates corrections that forecast the untreated treated--control gap and subtract its predicted post-minus-pre change from a common two-way fixed effects (TWFE) DID. A simulation compares no correction, linear and quadratic gap adjustments, Synthetic Control, Random Forest and Gradient Boosting across seven scenarios with a known treatment effect, 200 units, 20 periods and 500 replications per scenario. A polynomial correction matched to the data-generating process nearly eliminates signed bias, but improves root mean squared error (RMSE) only when the avoided bias exceeds the added forecast variance.

The same gap-correction identity is applied to balanced municipality/district-day homicide data for the 2020 police strike in Cear\'a, Brazil. The five empirical candidates are no correction, a linear group gap, Synthetic Control, Random Forest and Gradient Boosting. All share one 84-day pre-strike/13-day strike TWFE benchmark; 154 rolling pre-strike 84/13-day placebo windows select methods by RMSE without using strike-period treated outcomes. For total homicides, Random Forest has the lowest pre-strike RMSE (0.00809), narrowly ahead of unadjusted TWFE (0.00830). Its selected strike-period estimate is 0.03606 homicides per municipality/district-day versus TWFE of 0.03437. Random Forest is selected for six of seven outcomes, while no correction is selected for the fourth criminal-involvement indicator. Yet a 99-draw AIS bootstrap repeating the full selection selects no correction in 61 draws and Random Forest in only 15; its selected-procedure total-homicide interval $[-0.0254,0.1433]$ includes zero. With only five treated AIS, selection and causal inference remain fragile. Historical forecasting performance is not observed bias in the missing strike-period counterfactual. The joint evidence supports a conditional, auditable use of correction rather than a universally preferred forecast method.
\end{abstract}
